\chapter{Theory}
\label{chap:theory}
\section{Nonlinear Residual Problems}

Let 
\[
r : \mathbb{R}^n \to \mathbb{R}^m
\]
be a nonlinear and continuously differentiable function. The nonlinear residual problem is to find any
\begin{equation*}
     x^\ast \in \mathbb{R}^n \quad \text{such that} \quad r(x^*) = 0.
\end{equation*}
The function $r(x)$ is referred to as the \emph{residual} function.

\subsection{The Jacobian Matrix}

The Jacobian of \(r\) at a point \(x \in \mathbb{R}^n\) is the matrix
\[
J(x) = \frac{\partial r(x)}{\partial x}
=
\begin{bmatrix}
\dfrac{\partial r_1}{\partial x_1} & \cdots & \dfrac{\partial r_1}{\partial x_n} \\
\vdots & \ddots & \vdots \\
\dfrac{\partial r_m}{\partial x_1} & \cdots & \dfrac{\partial r_m}{\partial x_n}
\end{bmatrix}
\in \mathbb{R}^{m \times n}.
\]

Using first-order Taylor expansion, the Jacobian represents the first-order linearization of \(r\) near \(x\):
\begin{equation}
\label{eq:jac_res_approx}    
r(x + \Delta x) \approx r(x) + J(x)\,\Delta x.
\end{equation}

Thus, \(J(x)\) describes the local sensitivity of the residual with respect to perturbations in the variables.

\subsection{Existence and Number of Local Solutions}
\label{sec:sol_existence}

Let $x^\ast \in \mathbb{R}^n$ satisfy $r(x^\ast)=0$, and assume that $r$ is continuously differentiable in a neighborhood of $x^\ast$.
Then it follows from the Implicit Function Theorem that the local structure of solutions is determined by the rank of the Jacobian $J(x^\ast)$, provided the rank is constant in a neighborhood of $x^\ast$.

\paragraph{Case $m = n$ (square system).}

Suppose $m=n$ and let $x^\ast$ satisfy $r(x^\ast)=0$.  
If the Jacobian $J(x^\ast)$ is nonsingular,
\[
\det J(x^\ast) \neq 0,
\]
then, by the Implicit Function Theorem, $x^\ast$ is a locally unique solution. In particular, there exists a neighborhood around $x^\ast$ in which no other solutions exist.

If $J(x^\ast)$ is singular, the local behavior of the solution set cannot be determined from first-order information alone. In this case, the solution may fail to be isolated, and additional solutions may exist arbitrarily close to $x^\ast$.

\paragraph{Case $m < n$ (underdetermined system).}

If $J(x^\ast)$ has full row rank $m$, then the solution set near $x^\ast$ is not isolated. Instead, it forms a smooth manifold of dimension $n - m$.
Hence, there are infinitely many nearby solutions, parameterized by $n-m$ degrees of freedom.

\paragraph{Case $m > n$ (overdetermined system).}

In an overdetermined system, the number of equations exceeds the number of unknowns, so the system generally has no exact solution.
Suppose, however, that a point $x^\ast$ satisfies $r(x^\ast)=0$.
If $J(x^\ast)$ has full rank $n$, the constraints are locally independent and restrict all directions in $\mathbb{R}^n$, so the solution $x^\ast$ is locally isolated.
If $\operatorname{rank} J(x^\ast)<n$, the constraints are not independent, and the solution is structurally unstable under perturbations of $r$.

\paragraph{Global Solution.}

The Jacobian $J(x^\ast)$ characterizes only the \emph{local} structure of the solution set. Even if $J(x^\ast)$ is nonsingular and $x^\ast$ is locally unique, the nonlinear system may possess additional solutions elsewhere in $\mathbb{R}^n$. Consequently, local invertibility of the Jacobian does not imply global uniqueness.

Global existence or uniqueness of a solution to $r(x)=0$ requires additional structural assumptions on $r$, such as global injectivity or strong monotonicity. In the absence of such global properties, the existence and uniqueness of a solution cannot be guaranteed from an analysis of the Jacobian.

\section{Local Sensitivity of the Residual}

The Jacobian matrix characterizes how perturbations in the variables influence the residual locally (\cref{eq:jac_res_approx}). The magnitude of this influence may vary significantly depending on the direction of perturbation. Understanding these directional effects is important for assessing how reliably the variables can be inferred from the residual.\\

Local sensitivity can be analyzed through the singular values of the Jacobian, which quantify how strongly perturbations in different directions affect the residual.

\subsection{Singular Value Decomposition}

Let $J \in \mathbb{R}^{m \times n}$ denote the Jacobian evaluated at a point of interest. The singular value decomposition (SVD) of $J$ is given by
\begin{equation*}
J = U \Sigma V^\top,
\end{equation*}
where

\begin{itemize}
\item $U \in \mathbb{R}^{m \times m}$ and $V \in \mathbb{R}^{n \times n}$ are orthogonal matrices,
\item $\Sigma \in \mathbb{R}^{m \times n}$ is diagonal,
\item the diagonal entries satisfy
\[
\sigma_1 \ge \sigma_2 \ge \dots \ge \sigma_p \ge 0,
\]
with $p = \min(m,n)$.
\end{itemize}

The singular values describe how perturbations in specific directions of the variable space influence the residual. The columns of $V$ define orthogonal directions in $\mathbb{R}^n$, while the corresponding singular values determine the magnitude of the resulting change in the residual.\\

Large singular values correspond to directions in which small perturbations of the variables produce relatively large changes in the residual. Conversely, small singular values correspond to directions in which the residual changes only weakly. If one or more singular values are close to zero, the Jacobian is nearly rank-deficient, indicating directions along which the residual provides little local information about the variables.

\subsection{Condition Number of the Jacobian}

A scalar measure of local sensitivity is given by the condition number of the Jacobian,
\begin{equation}
    \label{eq:cond_num_def}
\kappa(J) = \frac{\sigma_1}{\sigma_p},
\end{equation}
where $\sigma_1$ and $\sigma_p$ denote the largest and smallest singular values of $J$, respectively.\\

The condition number measures how differently perturbations in various directions affect the residual. If the singular values are of comparable magnitude, perturbations in different directions produce changes in the residual of similar size, indicating balanced local sensitivity.\\

If the smallest singular value is much smaller than the largest, some directions produce much weaker changes in the residual than others. In such situations, small perturbations in the residual may correspond to large changes in the variables, indicating reduced numerical stability.\\

If $\sigma_p = 0$, the Jacobian is rank-deficient and the condition number is infinite. In this case, there exist directions in which perturbations of the variables do not influence the residual to first order.

\section{Optimization Methods}
\label{sec:theory_optmet}

In practice, it is often hard to solve a nonlinear residual problem directly. Instead, the problem can be reformulated as a nonlinear least squares optimization problem with the scalar objective function
\begin{equation}
f(x) = \frac{1}{2} \| r(x) \|^2 \in \mathbb{R}.
\label{eq:least_squares_objective}
\end{equation}
Solving the minimization problem,
\begin{equation}
\label{eq:min_problem}
    x^\ast = \underset{x}{\text{argmin}} \ f(x) = \underset{x}{\text{argmin}} \ \frac{1}{2} \| r(x) \|^2,
\end{equation}
satisfies $r(x^*) = 0$ if an exact solution exists. If no exact solution exists, or if the system is overdetermined, the formulation yields a solution that minimizes the residual norm.\\

Solving \cref{eq:min_problem} is very hard to solve in general \cite{madsen2004}, since there might exist local minima that are larger than a global minima. \cite{nielsen1999}

A group of methods used to solve the optimization problem operates by iteratively evaluating the objective function at the current iterate $x_k$ and determining a step $s_k$ that decreases the cost function. The next iterate is computed as
\begin{equation}
\label{eq:step_update_rule}
    x_{k+1} = x_k + s_k,
\end{equation}
with the goal of choosing $s_k$ so that
\begin{equation*}
    f(x_{k+1}) < f(x_k).
\end{equation*}
There are two main frameworks for controlling how the step $s_k$ is chosen, called \textit{Line Search} and \textit{Trust Region} methods. This section will introduce these frameworks along with common algorithms used to solve the minimization problem.

\subsection{Line Search Methods}
\begin{equation}
    s_k = \alpha_k p_k
\end{equation}
\begin{itemize}
    \item Armijo Backtracking\\
    \item Wolfe Conditions
\end{itemize}

\subsection{Gradient Descent}
Gradient descent is a simple method that is based on moving in the direction of the negative gradient of the objective function. 
Using the chain rule, the gradient of $f$ is given by
\begin{equation}
\label{eq:gradient_f}
    \nabla f(x) = J(x)^\top r(x),
\end{equation}
where $J(x)$ denotes the Jacobian of the residual function.

At the current iterate $x_k$, the search direction is chosen as the direction of steepest descent:
\begin{equation}
\label{eq:steepest_descent}
    p_k = - \nabla f(x_k) = - J(x_k)^\top r(x_k).
\end{equation}

Gradient descent is simple to implement and requires only first-order derivative information through the Jacobian. It is computationally inexpensive per iteration and applicable even when second-order derivatives are unavailable. However, the method may converge slowly, especially for ill-conditioned problems or narrow curved valleys in the objective function. Its performance is sensitive to scaling and may require many iterations to achieve high accuracy.

\subsection{Newton's Method}

Newton's method incorporates second-order information of the objective function. 
The gradient of $f$ is given by \cref{eq:gradient_f} and the Hessian is given by
\begin{equation}
\label{eq:hessian}
    \nabla^2 f(x) = J(x)^\top J(x) + \sum_{i=1}^m r_i(x)\nabla^2 r_i(x).
\end{equation}

At iteration $k$, the Newton search direction is obtained by solving the linear system
\begin{equation*}
    \nabla^2 f(x_k) p_k = - \nabla f(x_k).
\end{equation*}

By incorporating curvature information, Newton's method typically achieves quadratic convergence in a neighborhood of a solution, provided that the Hessian is nonsingular. This represents a significant improvement over gradient descent in terms of local convergence speed.

However, the method requires computation and factorization of the Hessian matrix, which can be expensive for large-scale problems. Moreover, if the Hessian is indefinite or poorly conditioned, the method may fail to produce a descent direction without additional safeguards.

\subsection{Gauss--Newton Method}

The Gauss--Newton method is tailored specifically to nonlinear least squares problems. 
It is based on an approximation of \cref{eq:hessian},
\begin{equation*}
    \nabla^2 f(x) \simeq J(x)^\top J(x),
\end{equation*}
which neglects the second-order term involving $\nabla^2 r_i(x)$.

The search direction is therefore computed by solving
\begin{equation}
\label{eq:gauss_newton}
    J(x_k)^\top J(x_k) p_k = - J(x_k)^\top r(x_k).
\end{equation}

This approximation avoids explicit computation of second derivatives and reduces computational cost compared to full Newton's method. When the residual is small near the solution, the neglected term becomes insignificant, and the method exhibits fast local convergence.

The Gauss--Newton method may perform poorly when the residual is large or when $J(x_k)^\top J(x_k)$ is rank deficient. In such cases, the approximation of the Hessian can be inaccurate, and additional regularization techniques may be required to ensure robustness.

\subsection{Trust Region Methods}
\label{sec:trust_region_methods}

Trust region methods determine the step $s_k$\footnote{The step $s_k$ is the same as the search direction $p_k$ with $\alpha=1$.} by minimizing a local model of the objective function within a restricted neighborhood around the current iterate. From \cite{madsen2004}, a typical model is the second order taylor expansion, using the Gauss--Newton approximation of the Hessian, 
\begin{equation}
\label{eq:mk}
f(x_k+s) \simeq m_k(s) = f(x_k) + r(x_k)^\top J(x_k) \, s + \frac{1}{2} s^\top J(x_k)^\top J(x_k) \, s
\end{equation}

In the trust region method, it is assumed that there exists a radius $\Delta$ around $x_k$, such that the model $m_k(s)$ is sufficiently accurate within the region. The step is then computed by solving the constrained minimization problem,

\begin{equation}
\label{eq:trust_region_subproblem}
    s_k = \underset{s}{\text{argmin}} \ m_k(s) \quad \text{subject to} \quad \|s\| \leq \Delta_k.
\end{equation}

In practice, this is solved through solving a damped problem, where the step is computed by solving the unconstrained minimization problem,
\begin{equation}
\label{eq:damped_problem}
    s_k = \underset{s}{\text{argmin}} \ \left(m_k(s) + \frac{1}{2} \ \lambda \ s^\top s\right),
\end{equation}
where $\lambda \geq 0$ is the damping parameter, and $\lambda \ s^\top s$ penalizes large steps. 

For a given $\Delta$ or $\lambda$, if no sufficient step $s_k$ is found such that $f(x_k+s_k)<f(x_k)$ (assuming $x_k$ is not a solution), the parameters ($\Delta,\lambda$) must be tuned such that a sufficient step is found. Since $m_k(s)$ is assumed to be a sufficient approximation for $f(x_k+s)$ when $s$ is sufficiently small, the trust region should be decreased ($\lambda$ increased) for the next iteration. When a sufficient step is found, it indicates that $m_k(s)$ is a good approximation, and the trust region can be increased ($\lambda$ decreased), aiming for faster convergence.

\subsection{Levenberg--Marquardt Method} \label{sec:levenberg_marquardt}

The Levenberg--Marquardt (LM) method is a damped Gauss--Newton method designed for solving nonlinear least--squares problems. It can be interpreted as a trust region method where the step is computed from a regularized Gauss--Newton system. The damping parameter, $\lambda$, controls the transition between the Gauss--Newton and steepest descent method.

At iteration $k$, the LM step is obtained by solving the linear system
\begin{equation} \label{eq:lm_system}
    \left( J(x_k)^\top J(x_k) + \lambda I \right) s_k
    =
    - J(x_k)^\top r(x_k),
\end{equation}
As $\lambda \to 0$, the system approaches the Gauss--Newton equations, \cref{eq:gauss_newton}. For large $\lambda$, the step approaches,
\begin{equation*}
    s_k \simeq -\frac{1}{\lambda}J(x_k)^\top r(x_k),
\end{equation*}
which is the steepest descent direciton, \cref{eq:steepest_descent}, with a small step lenght, $\alpha=\lambda^{-1}$.

The damping parameter must be initialized before the iterations begin. \cite{nielsen1999} states that the initial value should relate to the magnitude of the eigenvalues, and proposes
\begin{equation} \label{eq:init_lambda}
\lambda_0 = \tau \max \left( \operatorname{diag}(J(x_0)^\top J(x_0)) \right),
\end{equation}
where $\tau > 0$ is a small constant.

To evaluate the quality of a proposed step $s_k$, the LM method uses the gain ratio, defined as the ratio between the actual reduction in the objective function and the reduction predicted by the quadratic model,
\begin{equation}
\label{eq:gain_ratio}
\varrho =
\frac{f(x_k) - f(x_k + s_k)}
     {m_k(0) - m_k(s_k)}.
\end{equation}

The denominator can be computed from \cref{eq:mk},
\begin{align}
    m_k(0)-m_k(s_k) &= -s_k^\top J(x_k)^\top r(x_k) \,  - \frac{1}{2} s_k^\top J(x_k)^\top J(x_k) \, s_k \notag \\
    &= -\frac{1}{2} s_k^\top(2 J(x_k)^\top r(x_k) + (J(x_k)^\top J(x_k) + \lambda I-\lambda I)s_k) \notag \\
    &= \frac{1}{2} s_k^\top (\lambda s_k-J(x_k)^\top r(x_k)) \label{eq:mk_solve}.
\end{align}

Combining \cref{eq:gain_ratio,eq:mk_solve} yields a computationally easier expression for the gain ratio,
\begin{equation}
    \label{eq:gain_ratio_solve}
    \varrho = \frac{\|r(x_k)\|^2 - \|r(x_k+s_k)\|^2}
    {s_k^\top (\lambda s_k-J(x_k)^\top r(x_k))}.
\end{equation}
Note that $s_k^\top s_k$ and $-s_k^\top J(x_k)^\top r(x_k)$ are positive, such that the denominator is always positive.

If $\varrho > 0$, the step is accepted and the iterate is updated,
\[
x_{k+1} = x_k + s_k.
\]
The damping parameter is then decreased according to
\begin{equation}
\label{eq:dec_lambda}
\lambda := 
\lambda \max \left\{ \frac{1}{3},\, 1 - (2\varrho - 1)^3 \right\}.
\end{equation}

If $\varrho \le 0$, the step is rejected, and the damping parameter is increased by
\begin{equation}
\label{eq:inc_lambda}
\lambda := 2^{i+1}\lambda,
\end{equation}
where $i\in \{0,1,\dots,i_{\max}\}$ is the iteration counter until the step is accepted. This update strategy is demonstrated in \cite{nielsen1999}, and is shown to be a sufficient choice.\\

Intuitively, when $\varrho$ is positive, the quadratic model $m_k(s)$ provides a good approximation of the objective function, and the damping parameter is reduced so that the method behaves more like Gauss--Newton. When $\varrho$ is negative, the model is considered unreliable, and the damping parameter is increased, producing a smaller and more conservative step.


\subsubsection{Computing the step with SVD factorization}
\cref{eq:lm_system} can be solved by SVD factorization. Using,
\begin{equation*}
    J(x_k) = U_k \Sigma_k V_k^\top,
\end{equation*}
and performing some algebraic manipulation, the step can be computed as
\begin{equation}
\label{eq:sk_lm_solve}
    s_k = -V_k \, \text{diag}(d_i) \, U_k^\top r(x_k), 
\end{equation}
Where $d_i=\dfrac{\sigma_i}{\sigma_i^2 + \lambda}$, and $\sigma_i$ are the singular values from $\Sigma_k$.\\

\cref{eq:sk_lm_solve} computes the LM step without computing $J^\top J$, which can be numerically unstable when $J$ is ill-conditioned.









